\documentclass[12pt]{article}
\usepackage{amsmath, amssymb, geometry}
\geometry{margin=1in}
\setlength{\parindent}{0pt}
\setlength{\parskip}{0.6em}

\title{Practice Problem Set 3}
\author{ECON 300 -- Labour Economics}
\date{Winter 2026}

\begin{document}
\maketitle


\textbf{1. Competitive Labour Market}

Suppose the market labour supply is:
\[
N^S = 20W
\]

and market labour demand is:
\[
N^D = 400 - 10W
\]

\begin{enumerate}
\item[(a)] Find the competitive equilibrium wage and employment.
\item[(b)] Compute worker surplus and firm surplus.
\item[(c)] Graph the market and clearly label the equilibrium.
\item[(d)] Suppose labour demand increases to:
\[
N^D = 500 - 10W
\]
Find the new equilibrium and explain the adjustment mechanism.
\end{enumerate}


\vspace{1em}

\textbf{2. Payroll Tax Incidence}

Using the original market in Question 1, suppose a payroll tax of $T = 5$ per worker is imposed on employers.

\begin{enumerate}
\item[(a)] Write the new labour demand equation in terms of the wage received by workers.
\item[(b)] Find the new wage paid by firms and the wage received by workers.
\item[(c)] Compute the new level of employment.
\item[(d)] Calculate the tax burden borne by workers and firms.
\item[(e)] Derive the deadweight loss algebraically.
\end{enumerate}


\vspace{1em}

\textbf{3. Labour Demand: Competitive vs Monopolist Output Market}

A competitive firm has:
\[
MPP_N = 50 - 2N
\]
Output price is $P = 10$.

\begin{enumerate}
\item[(a)] Derive the firm’s labour demand curve.
\item[(b)] If the wage is 200, find optimal employment.
\end{enumerate}

Now suppose the firm is a monopolist with:
\[
P = 20 - Q
\]
and production function:
\[
Q = 10N
\]

\begin{enumerate}
\item[(c)] Derive marginal revenue.
\item[(d)] Derive the marginal revenue product of labour.
\item[(e)] Compare employment under monopoly versus competition at wage 200.
\end{enumerate}


\vspace{1em}

\textbf{4. Basic Monopsony Model}

Suppose a monopsonist faces labour supply:
\[
W = 10 + 0.5N
\]

\begin{enumerate}
\item[(a)] Derive the total labour cost function.
\item[(b)] Derive the marginal cost of labour.
\item[(c)] Suppose:
\[
VMP_N = 100 - N
\]
Find the profit-maximizing employment and wage.
\item[(d)] Compare this outcome to the competitive equilibrium.
\item[(e)] Compute the monopsony wage gap.
\end{enumerate}


\vspace{1em}

\newpage
\textbf{5. Minimum Wage under Monopsony}

Using Question 4, suppose the government sets a minimum wage of 50.

\begin{enumerate}
\item[(a)] Determine whether employment increases or decreases.
\item[(b)] Derive the new marginal cost of labour under the binding minimum wage.
\item[(c)] Determine employment.
\item[(d)] Is this minimum wage efficiency-enhancing? Explain.
\end{enumerate}


\vspace{1em}

\textbf{6. Perfect Monopsonistic Wage Discrimination}

Using the labour supply and $VMP_N$ from Question 4:

\begin{enumerate}
\item[(a)] Derive employment under perfect wage discrimination.
\item[(b)] Compare employment under monopsony, perfect discrimination, and competition.
\item[(c)] Compute total worker surplus under discrimination.
\item[(d)] Explain why employment is efficient but distribution differs.
\end{enumerate}


\vspace{1em}

\textbf{7. Short Run vs Long Run Labour Supply to the Firm}

Suppose a firm faces short-run supply:
\[
W = 20 + 0.2N
\]
and long-run perfectly elastic supply at:
\[
W = 40
\]
Labour demand is:
\[
VMP_N = 100 - 2N
\]

\begin{enumerate}
\item[(a)] Find short-run equilibrium wage and employment.
\item[(b)] Find long-run equilibrium employment.
\item[(c)] Explain why wage falls in the long run.
\item[(d)] Graph the adjustment path.
\end{enumerate}

\end{document}